% ==============================================================================
% core_step4_sysid.tex — SCI/SICE tutorial core deck, Step 4: system
% identification (exercise 7, `sf sysid fit` only)
% SCI/SICE チュートリアル要点版 Step 4: システム同定（実習 7。sf sysid fit のみ）
% ==============================================================================

% ------------------------------------------------------------------
\begin{frame}{システム同定: 飛行データから $K,\tau_m$ を推定する}
  \begin{columns}[c]
    \begin{column}{0.46\textwidth}
      \centering
      \adjustbox{max width=\columnwidth, max height=0.78\textheight}{\documentclass[border=10pt]{standalone}
\usepackage{tikz}
\usepackage{amsmath}
\usetikzlibrary{arrows.meta, positioning, shapes.geometric, calc}

\begin{document}

\definecolor{blockcolor}{RGB}{0, 120, 200}
\definecolor{arrowcolor}{RGB}{80, 80, 80}
\definecolor{plantcolor}{RGB}{40, 160, 40}
\definecolor{estcolor}{RGB}{200, 100, 0}
\begin{tikzpicture}[
    block/.style={
      rectangle, draw=#1, line width=1.5pt,
      fill=#1!8, rounded corners=3pt,
      minimum width=38mm, minimum height=16mm,
      font=\bfseries, text=#1, align=center
    },
    arrow/.style={-{Stealth[length=3mm]}, line width=1.2pt, arrowcolor},
    lbl/.style={font=\small, above, text=arrowcolor},
  ]

  % ---------- Top row: u(t) → Plant → y(t) ----------

  % Input
  \node[font=\bfseries, arrowcolor] (input) at (-1, 0) {$u(t)$};

  % Plant block (centered)
  \node[block=plantcolor] (plant) at (4.5, 0)
    {Plant $G_p(s)$\\[2pt]{\small (unknown $K$, $\tau_m$)}};

  % Output
  \node[font=\bfseries, arrowcolor] (output) at (10, 0) {$y(t)$};

  % Main signal flow
  \draw[arrow] (input) -- (plant);
  \draw[arrow] (plant) -- (output);

  % ---------- Bottom: Estimation (same center as Plant) ----------

  \node[block=estcolor] (est) at (4.5, -3.5)
    {Parameter\\Estimation};

  % ---------- Branch arrows: u → est.west, y → est.east ----------

  % Tap points on signal lines (symmetric about center)
  \coordinate (tap_u) at (1.0, 0);
  \coordinate (tap_y) at (8.0, 0);

  % Junction dots
  \fill[arrowcolor] (tap_u) circle (2.5pt);
  \fill[arrowcolor] (tap_y) circle (2.5pt);

  % u(t) branch: down then right → est.west
  \draw[arrow] (tap_u) |- (est.west);

  % y(t) branch: down then left → est.east
  \draw[arrow] (tap_y) |- (est.east);

  % ---------- Estimation output ----------

  \node[font=\bfseries, estcolor] (params) at (4.5, -5.5) {$\hat{K}$, $\hat{\tau}_m$};
  \draw[arrow] (est) -- (params);

  % ---------- Labels ----------

  \node[font=\scriptsize, arrowcolor, above=1mm] at (input.north)
    {control input};
  \node[font=\scriptsize, arrowcolor, above=1mm] at (output.north)
    {measured output};
  \node[font=\scriptsize, estcolor, below=2mm] at (params.south)
    {compare with L6 theory};

\end{tikzpicture}
\end{document}
}
      {\tiny 図中の L6 = 実習 6 の理論値\par}
    \end{column}
    \begin{column}{0.52\textwidth}
      \footnotesize
      \setlength{\abovedisplayskip}{3pt}\setlength{\belowdisplayskip}{3pt}%
      入力と出力の記録から、プラント $G_p(s)=K/(s(\tau_m s+1))$ の $K,\tau_m$ を求める。一般には、候補の $K,\tau_m$ でモデルの出力 $\hat y(t)$ を計算し、実測 $y(t)$ との二乗誤差が最小になる値を探す（最小二乗フィット）:
      \[
        J(K,\tau_m)=\sum_t\bigl(\hat y(t)-y(t)\bigr)^2
      \]
      \code{sf sysid fit} は既定で、閉ループ（P 制御で飛ばした記録）の「目標角速度 $\to$ ジャイロ」の応答にモデルを当てはめる。制御なしでは飛べない機体なので閉ループで同定する。$K_p$ は記録のモータ Duty から自動で推定する
    \end{column}
  \end{columns}
\end{frame}

% ------------------------------------------------------------------
\begin{frame}[fragile]{データの取り方: 機体の WiFi につなぐ}
  \begin{columns}[T]
    \begin{column}{0.50\textwidth}
      \begin{exampleblock}{PC を機体の WiFi につなぐ}
        \footnotesize
        \begin{itemize}\setlength{\itemsep}{1pt}
          \item 実習 1 の \code{wifi.mode 1} で、機体が自分の WiFi を出している
          \item SSID は \code{StampFly-XXYY}（XXYY = 機体の MAC アドレス末尾 4 桁。機体のラベルと同じ）、パスワードは \code{stampfly}
          \item 機体の IP は \code{192.168.10.1}。接続中、PC は会場のネットから切り替わる
        \end{itemize}
      \end{exampleblock}
    \end{column}
    \begin{column}{0.48\textwidth}
      \begin{table}
        \centering\scriptsize
        \begin{tabular}{>{\raggedright\arraybackslash}p{30mm}>{\raggedright\arraybackslash}p{32mm}}
          \hline
          \textbf{方法} & \textbf{使う場面} \\
          \hline
          Teleplot（USB） & その場で波形を見る（実習 2） \\
          \code{sf telemetry}（WiFi） & コード変更なしで 50\,Hz の状態を見る \\
          \code{sf log wifi}（WiFi） & 後で解析するため 400\,Hz で記録する \\
          \hline
        \end{tabular}
      \end{table}
      {\scriptsize WiFi は飛行中も使える（USB は飛行中に挿せない）。今日の実習 7 では \code{sf log wifi} を使う\par}
    \end{column}
  \end{columns}
\end{frame}

% ------------------------------------------------------------------
\begin{frame}[fragile]{実習 7: ログを解析し、$K,\tau_m$ を求める}
  \exgoal{実習 6 の理論値 $K,\tau_m$ が、実機の飛行データと一致するか確かめる}
  \begin{columns}[T]
    \begin{column}{0.55\textwidth}
      {\footnotesize\textbf{\textcolor{sfblue}{手順}}}\par\vspace{-2pt}
      \footnotesize
      \begin{enumerate}\setlength{\itemsep}{0pt}
        \item \code{sf lesson switch sci2026:7}、\code{sf lesson build}（TODO は \api{ws::set\_rate\_target} の呼び出し。目標角速度をログに記録する。実習 5 の雛形にも同じ TODO がある）
        \item \textbf{講師}が飛ばし、\code{sf log wifi -d 30} で記録する。ログは配布される
        \item \textbf{各自}: 配布されたログ（自分で取ったログでもよい）に \code{sf sysid fit}
      \end{enumerate}
    \end{column}
    \begin{column}{0.43\textwidth}
      \begin{sflisting}[basicstyle=\ttfamily\scriptsize]{コマンド}
sf log list
sf sysid fit <ログ名> --plot
      \end{sflisting}
      \vspace{-1mm}
      {\scriptsize \code{.sflog.zip} は信号ごとの CSV をまとめた zip（フライトログ一式）。\code{sf log wifi} は \code{logs/flight\_<日時>.sflog.zip} に保存し、\code{sf log list} で名前を確認できる\par}
    \end{column}
  \end{columns}
  \exresult{軸ごとの $K,\tau_m$ と当てはまりの良さ $R^2$ が表示され、実習 6 の理論値と比べられる}
  \exhelp{7}
\end{frame}

% ------------------------------------------------------------------
\begin{frame}{同定結果と理論値の比較（例）}
  \begin{table}
    \centering\footnotesize
    \begin{tabular}{lccccc}
      \hline
      \textbf{軸} & \textbf{$K$（同定）} & \textbf{$K$（理論）} & \textbf{$\tau_m$（同定）} & \textbf{$\tau_m$（理論）} & \textbf{$R^2$} \\
      \hline
      Roll  & 98.5 & 102.0 & 0.019\,s & 0.020\,s & 0.94 \\
      Pitch & 72.1 &  70.0 & 0.021\,s & 0.020\,s & 0.91 \\
      \hline
    \end{tabular}
  \end{table}
  {\footnotesize \textbf{上の数値は読み方を示すための例（説明用の数値）で、実測ではない}。実際の値は飛行ごとに変わる\par}

  \vspace{0.4em}

  \begin{itemize}\footnotesize\setlength{\itemsep}{3pt}
    \item $K$ が理論値の数 \% 以内で一致すれば、実習 6 のモデルは実機を説明できている
    \item $R^2$ は 1 に近いほど、モデルがデータをよく説明している
    \item 人がスティックで操縦した記録は励振（入力で系を十分に動かすこと）が小さいので、当てはまりが悪いときは別のログで試す
  \end{itemize}
\end{frame}
